\documentclass[10pt,a4paper]{extarticle}
\usepackage{parskip}

\usepackage[utf8]{inputenc}
\usepackage[T1]{fontenc}
\usepackage[brazil]{babel}
\usepackage{amsmath, amssymb}
\usepackage{geometry}
\usepackage{xcolor}
\usepackage{hyperref}
\usepackage{enumitem}
\usepackage{microtype}
\usepackage{csquotes}
\usepackage{natbib}

\definecolor{important}{RGB}{255, 100, 100}
\definecolor{notice}{RGB}{0, 102, 204}

\newcommand{\impt}[1]{\textcolor{important}{#1}}
\newcommand{\ntc}[1]{\textcolor{notice}{#1}}

\setlist{nosep}
\geometry{margin=2cm}

\hypersetup{
  colorlinks=true,
  linkcolor=notice,
  citecolor=notice,
  urlcolor=notice
}

\title{Gabarito -- Estudo Dirigido II -- Falsificacionismo\footnote{Respostas baseadas em \citep{chalmers1993}.}}
\author{BCC502 -- Metodologia Científica para Ciência da Computação}
\date{\today}

\begin{document}

\vspace{-1cm}
\maketitle

\section*{Capítulo 4}
\begin{enumerate}[label=\textbf{\arabic*.}, itemsep=10pt]

\item \textbf{Qual é a visão do falsificacionista a respeito das observações e teorias?}
\begin{itemize}
  \item Aceita que a observação é orientada e pressuposta pela teoria, e não a considera neutra ou anterior a ela.
  \item Rejeita a ideia de que teorias possam ser estabelecidas como verdadeiras ou provavelmente verdadeiras a partir da evidência observacional.
  \item Vê as teorias como conjecturas especulativas, criadas livremente pelo intelecto humano para superar problemas deixados por teorias anteriores.
  \item Exige que toda teoria proposta seja submetida a testes rigorosos e inexoráveis de observação e experimento.
  \item Teorias que falham nesses testes devem ser eliminadas e substituídas por novas conjecturas.
  \item Uma teoria nunca pode ser dita verdadeira; no máximo, pode-se dizer que é a melhor disponível até agora, por ter resistido mais do que as anteriores.
\end{itemize}

\item \textbf{Qual é a visão do falsificacionista em relação ao progresso da ciência?}
\begin{itemize}
  \item A ciência progride por tentativa e erro: por conjecturas e refutações.
  \item Apenas as teorias mais bem adaptadas, isto é, as que resistem aos testes, sobrevivem.
  \item Embora nunca se possa chamar uma teoria de verdadeira, pode-se dizer que a teoria vigente é melhor que qualquer outra proposta anteriormente.
  \item O crescimento do conhecimento científico é o conceito central dessa explicação da ciência, detalhado no capítulo seguinte.
\end{itemize}

\item \textbf{Como o falsificacionismo pode ser justificado logicamente?}
\begin{itemize}
  \item Baseia-se numa assimetria lógica entre verificação e falsificação.
  \item Não é logicamente possível derivar leis e teorias universais a partir de proposições de observação singulares (o problema da indução).
  \item É logicamente possível, porém, deduzir a falsidade de leis e teorias universais a partir de proposições de observação singulares verdadeiras.
  \item Exemplo do livro: da premissa \enquote{um corvo não preto foi observado no local $x$ no momento $n$} deduz-se validamente que \enquote{nem todos os corvos são pretos}.
  \item É essa assimetria, dedução válida da falsidade mas não da verdade, que o falsificacionista explora ao máximo.
\end{itemize}

\item \textbf{O que é uma hipótese falsificável?}
\begin{itemize}
  \item É uma hipótese para a qual existe ao menos uma proposição de observação logicamente possível que, se estabelecida como verdadeira, a contradiria.
  \item Não precisa ser falsa: basta que corra o risco lógico de sê-lo (exemplos do livro: \enquote{nunca chove às quartas-feiras}, \enquote{objetos pesados caem diretamente para baixo quando soltos}).
  \item Não são falsificáveis as tautologias (\enquote{ou está chovendo ou não está}), as verdades por definição (\enquote{todos os pontos de um círculo euclidiano são equidistantes do centro}) e as afirmações vagas o suficiente para se ajustarem a qualquer resultado (horóscopos, e a caricatura que Chalmers faz da psicologia adleriana).
  \item A falsificabilidade é exigida porque é o que dá conteúdo informativo a uma lei: uma afirmação infalsificável é compatível com qualquer comportamento do mundo e, por isso, nada informa sobre ele.
\end{itemize}

\item \textbf{Como se estabelece a preferência entre uma teoria e outra dentro do falsificacionismo?}
\begin{itemize}
  \item Prefere-se a teoria mais falsificável, isto é, a que faz afirmações mais amplas e arriscadas sobre o mundo, desde que ela ainda não tenha sido falsificada.
  \item Compara-se o conjunto de \enquote{falsificadores potenciais} de cada teoria (as proposições de observação que a falsificariam): se os falsificadores potenciais de uma teoria formam um subconjunto dos de outra, esta última é mais falsificável.
  \item Exemplos do livro: \enquote{todos os planetas se movem em elipses} é preferível a \enquote{Marte se move em elipse}; a teoria de Newton, com muito mais falsificadores potenciais que a de Kepler, é preferível a ela por ter resistido aos testes.
  \item Teorias falsificadas devem ser rejeitadas inexoravelmente; entre as que resistem, vence a mais falsificável.
\end{itemize}

\item \textbf{Qual é a consequência atrativa que vem da exigência de que as teorias sejam altamente falsificáveis?}
\begin{itemize}
  \item A exigência de alta falsificabilidade obriga as teorias a serem formuladas com clareza e precisão.
  \item Teorias vagas podem sempre ser reinterpretadas como compatíveis com qualquer resultado de teste, escapando assim à falsificação (exemplo do livro: a descrição vaga que Goethe faz da eletricidade).
  \item Quanto mais precisa a formulação, mais falsificável a teoria: \enquote{os planetas movem-se em elipses em torno do Sol} é mais falsificável, e por isso preferível, a \enquote{os planetas movem-se em curvas fechadas em torno do Sol}.
  \item Essa exigência impede manobras de proteção como as de políticos e videntes, que se resguardam com afirmações vagas demais para serem contrariadas por qualquer observação.
\end{itemize}

\item \textbf{Como se dá o progresso da ciência na visão do falsificacionista?}
\begin{itemize}
  \item A ciência começa com problemas, não com observação neutra: um fenômeno só é problemático à luz de alguma teoria já aceita.
  \item Cientistas propõem hipóteses falsificáveis como soluções tentativas para o problema; essas hipóteses são criticadas e testadas.
  \item Algumas hipóteses são rapidamente eliminadas; outras resistem e passam a ser submetidas a críticas e testes ainda mais rigorosos.
  \item Quando uma hipótese bem testada é eventualmente falsificada, surge um novo problema, e o ciclo recomeça com novas conjecturas.
  \item Exemplo simples do livro: o voo dos morcegos, com hipóteses sucessivas sendo testadas e falsificadas (visão, orelhas, ouvidos) até chegar à hipótese da ecolocalização pelos guinchos.
  \item Exemplo mais amplo: a física de Aristóteles, falsificada e substituída pela de Newton, que por sua vez é falsificada (órbita de Mercúrio, massa dos elétrons) e substituída pela de Einstein.
\end{itemize}

\end{enumerate}

\section*{Capítulo 5}
\begin{enumerate}[resume, label=\textbf{\arabic*.}, itemsep=10pt]

\item \textbf{Quais condições uma hipótese deve satisfazer?}
\begin{itemize}
  \item Deve ser falsificável: deve haver proposições de observação logicamente possíveis que a contradigam.
  \item Quanto mais falsificável, melhor, pois maior é seu conteúdo informativo.
  \item Não deve, até o momento, ter sido falsificada pelas provas disponíveis.
  \item Para o falsificacionista mais sofisticado, ainda é preciso que a hipótese seja mais falsificável do que a teoria que pretende substituir, de modo a garantir o progresso da ciência.
\end{itemize}

\item \textbf{Qual é a explicação falsificacionista sofisticada da ciência?}
\begin{itemize}
  \item Desloca o foco dos méritos de uma teoria isolada para os méritos relativos entre teorias concorrentes: uma visão dinâmica, e não estática, da ciência.
  \item A pergunta central deixa de ser \enquote{esta teoria é falsificável? quão falsificável? foi falsificada?} e passa a ser \enquote{esta nova teoria é um substituto viável para a teoria contestada?}
  \item Uma teoria nova é aceita como digna de consideração se for mais falsificável que sua rival, sobretudo se prevê um novo tipo de fenômeno não tocado por ela.
  \item Reconhece também um papel importante para a confirmação, e não apenas para a falsificação: a ciência avança tanto pela confirmação de conjecturas audaciosas quanto pela falsificação de conjecturas cautelosas.
\end{itemize}

\item \textbf{Qual é o problema técnico do falsificacionismo ingênuo? Como o falsificacionismo sofisticado o resolve?}
\begin{itemize}
  \item O problema: não é possível definir uma medida absoluta do grau de falsificabilidade de uma teoria isolada, pois o número de seus falsificadores potenciais é sempre infinito.
  \item Não há, portanto, como responder objetivamente a uma pergunta como \enquote{quão falsificável é a lei da gravidade de Newton?}.
  \item O falsificacionismo sofisticado resolve o problema abandonando a ideia de grau absoluto e adotando graus relativos de falsificabilidade.
  \item Compara-se o conjunto de falsificadores potenciais de duas teorias: se os de uma formam um subconjunto dos de outra, esta última é mais falsificável.
  \item A evolução histórica de uma ciência é então avaliada por essa comparação sucessiva, cada teoria da série sendo, idealmente, mais falsificável que sua predecessora.
\end{itemize}

\item \textbf{O que é uma modificação \emph{ad hoc}? Apresente um exemplo de uma modificação \emph{ad hoc} e um de uma modificação não \emph{ad hoc}.}
\begin{itemize}
  \item Modificação \emph{ad hoc}: alteração numa teoria (acréscimo ou mudança de postulado) que não leva a nenhuma consequência testável nova além das que a teoria já tinha; serve apenas para proteger a teoria de uma falsificação que a ameaça.
  \item Exemplo de modificação \emph{ad hoc}: diante do envenenamento causado por um lote de pão numa aldeia francesa, restringir \enquote{todo pão alimenta} para \enquote{todo pão, exceto aquele lote específico, alimenta}, o que não abre nenhum teste além do próprio consumo do pão.
  \item Outro exemplo do livro: a resposta do adversário aristotélico de Galileu à observação de crateras na Lua, postulando uma substância invisível e indetectável que preencheria as crateras e cobriria as montanhas.
  \item Exemplo de modificação não \emph{ad hoc}: reformular a hipótese para \enquote{todo pão alimenta, exceto o feito de trigo contaminado por um fungo específico}, o que abre novos testes independentes (análise química do fungo, cultivo em trigo preparado, etc.).
  \item Outro exemplo do livro: a hipótese de Leverrier e Adams sobre um planeta ainda não detectado perturbando a órbita de Urano, que levou a um novo teste independente (a observação telescópica de Galle) e culminou na descoberta de Netuno.
\end{itemize}

\item \textbf{De onde advêm as principais contribuições ao crescimento do conhecimento científico? Por quê?}
\begin{itemize}
  \item Advêm da confirmação de conjecturas audaciosas ou da falsificação de conjecturas cautelosas.
  \item Ambos os casos são informativos: o primeiro assinala a descoberta de algo previamente desconhecido ou tido como improvável; o segundo revela que algo tido como certo é, na verdade, falso.
  \item Exemplos de confirmação de conjectura audaciosa: a descoberta de Netuno, a detecção das ondas de rádio e a confirmação de Eddington da curvatura da luz prevista por Einstein.
  \item Em contraste, a falsificação de uma conjectura audaciosa (a especulação de Kepler sobre os sólidos platônicos) ou a confirmação de uma conjectura cautelosa (mais uma amostra de aço que se expande ao ser aquecida) pouco acrescentam ao conhecimento científico.
\end{itemize}

\item \textbf{Qual é a meta da ciência no falsificacionismo?}
\begin{itemize}
  \item Falsificar teorias vigentes e substituí-las por teorias melhores, isto é, mais falsificáveis, com maior conteúdo informativo e testável.
  \item As confirmações de teorias novas importam na medida em que demonstram que a nova teoria é um aperfeiçoamento em relação à teoria que está sendo substituída.
  \item Uma teoria audaciosa recém-confirmada torna-se, por sua vez, o novo alvo a ser testado por conjecturas ulteriores ainda mais audaciosas, num processo sem fim.
\end{itemize}

\item \textbf{Como se diferem as visões indutivista e falsificacionista sobre o papel da observação?}
\begin{itemize}
  \item Para o indutivista, o valor de uma instância confirmadora depende só da relação lógica entre a proposição de observação e a teoria, independentemente do contexto histórico; quanto mais confirmações, maior o apoio indutivo e a probabilidade de a teoria ser verdadeira.
  \item Para o falsificacionista sofisticado, o valor de uma confirmação é essencialmente histórico: só é significativa se resulta do teste de uma previsão nova, improvável à luz do conhecimento prévio da época.
  \item Confirmações que são conclusões inevitáveis à luz do conhecimento já estabelecido são cientificamente insignificantes para o falsificacionista, embora sejam logicamente equivalentes a confirmações significativas.
  \item Exemplo do livro: a detecção das ondas de rádio por Hertz foi uma confirmação historicamente relevante da teoria de Maxwell; ouvir rádio hoje é, logicamente, uma confirmação idêntica, mas sem nenhum valor científico.
\end{itemize}

\end{enumerate}

\section*{Capítulo 6}
\begin{enumerate}[resume, label=\textbf{\arabic*.}, itemsep=10pt]

\item \textbf{Como se diferem o \enquote{falsificacionista ingênuo} e o \enquote{falsificacionista mais sofisticado}?}
\begin{itemize}
  \item O falsificacionista ingênuo concentra-se apenas nas tentativas de falsificar teorias, estabelecendo a verdade de proposições de observação inconsistentes com elas, e avalia cada teoria isoladamente perguntando apenas se ela é falsificável e se foi falsificada.
  \item O falsificacionista sofisticado reconhece também a importância da confirmação de teorias especulativas e compara graus relativos de falsificabilidade entre teorias concorrentes, adotando uma visão dinâmica do progresso científico.
  \item Ambos compartilham, porém, a suposição de que há uma diferença qualitativa entre confirmação (sempre provisória) e falsificação (potencialmente decisiva), suposição que Chalmers questiona neste capítulo.
\end{itemize}

\item \textbf{Justifique a afirmação a seguir: \enquote{Todas as proposições são falíveis}.}
\begin{itemize}
  \item Retoma o argumento do Capítulo III: a observação depende da teoria, e não existe uma base observacional pura e infalível.
  \item Como toda proposição de observação pressupõe alguma teoria, ela pode ser posta em dúvida por desenvolvimentos teóricos posteriores.
  \item Nada na lógica da situação exige que, diante de um choque entre teoria e observação, seja sempre a teoria (e não a proposição de observação) a ser rejeitada.
  \item Exemplo do livro: a observação a olho nu de que Vênus não muda de tamanho apreciavelmente ao longo do ano foi tratada, a princípio, como inconsistente com a teoria de Copérnico; um século depois, essa proposição de observação foi revista e rejeitada graças a avanços na óptica, e a teoria de Copérnico foi mantida.
  \item Consequência: se as proposições de observação são falíveis, falsificações conclusivas de teorias não são de fato alcançáveis.
\end{itemize}

\item \textbf{Quais são as objeções feitas ao falsificacionismo?}
\begin{itemize}
  \item A observação depende da teoria e é falível, de modo que falsificações nunca podem ser conclusivas.
  \item A defesa que Popper oferece contra essa dificuldade não a resolve, pois as proposições de observação continuam falíveis (ver questão seguinte).
  \item Situações reais de teste são complexas: uma teoria só é testada em conjunto com hipóteses auxiliares e condições iniciais, de modo que uma previsão falsa não indica sozinha qual premissa está errada.
  \item O falsificacionismo é inadequado em bases históricas: teorias hoje consideradas exemplares da ciência foram mantidas por seus autores mesmo depois de terem sido, a rigor, falsificadas logo no início.
  \item A Revolução Copernicana não pode ser explicada satisfatoriamente nem pelo indutivismo nem pelo falsificacionismo, dada a complexidade e a duração da mudança teórica envolvida.
\end{itemize}

\item \textbf{Como Popper defende o falsificacionismo dessas objeções? Por que essa defesa não salva o falsificacionismo? Dê exemplos.}
\begin{itemize}
  \item Popper distingue as proposições de observação públicas (as \enquote{afirmações básicas}, testáveis e formuladas em linguagem pública) das experiências perceptivas privadas dos observadores individuais, que apenas motivam, mas não estabelecem, a aceitação de uma proposição.
  \item Uma afirmação básica é aceita ou rejeitada por decisão da comunidade científica, com base em sua capacidade de sobreviver a testes (ver também a questão seguinte).
  \item Popper reconhece, com a metáfora das estacas fincadas num pântano em vez de sobre uma base rochosa sólida, que a base empírica da ciência não é absoluta.
  \item Essa defesa não salva o falsificacionismo porque, mesmo assim, as proposições de observação continuam falíveis e sujeitas a revisão: uma falsificação aceita hoje pode ser desfeita amanhã por novos desenvolvimentos teóricos, de modo que falsificações conclusivas permanecem inatingíveis.
  \item Exemplo do livro: a observação a olho nu do tamanho aparente constante de Vênus, tida como inconsistente com a teoria de Copérnico na época, foi revogada cerca de cem anos depois, graças a desenvolvimentos na óptica.
\end{itemize}

\item \textbf{Quando uma proposição de afirmação é aceitável?}
\begin{itemize}
  \item Na reformulação que Chalmers propõe para a posição de Popper, uma proposição de observação é aceitável, experimentalmente, se, num dado estágio do desenvolvimento de uma ciência, ela é capaz de passar por todos os testes tornados possíveis pelo estado de desenvolvimento da ciência naquele estágio.
  \item Não depende, portanto, de uma percepção sensorial infalível de um indivíduo, mas de sua capacidade de resistir à crítica pública e aos testes disponíveis no momento.
  \item Por isso, sua aceitação é sempre provisória: pode ser rejeitada mais tarde à luz de testes futuros, possibilitados por novos desenvolvimentos teóricos e técnicos.
\end{itemize}

\item \textbf{Por que a aplicação do falsificacionismo em situações de teste realistas é complexa? Dê exemplos.}
\begin{itemize}
  \item Uma teoria científica real não é uma afirmação isolada, mas um complexo de leis, acrescido de suposições auxiliares (por exemplo, as teorias sobre o funcionamento dos instrumentos usados) e de condições iniciais (a descrição do cenário experimental).
  \item Quando uma previsão derivada desse conjunto de premissas falha, a lógica só garante que ao menos uma das premissas é falsa, sem indicar qual: pode ser a teoria em teste, uma suposição auxiliar ou as condições iniciais.
  \item Exemplo: o desvio da órbita de Urano em relação à prevista pela teoria de Newton não indicava um erro na teoria, mas uma condição inicial incompleta, a existência, ainda desconhecida, do planeta Netuno.
  \item Exemplo: Tycho Brahe julgou ter refutado a teoria copernicana por não detectar a paralaxe estelar prevista; na verdade, o erro estava em sua suposição sobre a distância das estrelas fixas, não na teoria de Copérnico.
  \item Exemplo hipotético de Lakatos, citado por Chalmers: um físico defende a teoria de Newton de sucessivas aparentes refutações propondo, a cada vez, uma nova hipótese auxiliar (um planeta perturbador, depois uma nuvem de poeira cósmica, depois um campo magnético), mostrando como uma teoria pode sempre ser protegida desviando-se a falsificação para outra parte do conjunto de suposições.
\end{itemize}

\item \textbf{Quais são as evidências da inadequação histórica do falsificacionismo?}
\begin{itemize}
  \item Teorias hoje consideradas exemplares da ciência foram, a rigor, falsificadas logo que propostas, e mesmo assim os cientistas não as abandonaram, o que foi bom para a ciência.
  \item A teoria gravitacional de Newton foi falsificada nos primeiros anos por observações da órbita lunar (a falsificação só foi atribuída a outra causa quase cinquenta anos depois) e, mais tarde, pelos detalhes da órbita de Mercúrio, sem nunca ter sido protegida de modo satisfatório.
  \item A teoria atômica de Bohr era, em suas versões iniciais, inconsistente com a observação de que a matéria é estável por mais de cerca de $10^{-8}$ segundos.
  \item A teoria cinética dos gases de Maxwell foi reconhecida como falsificada pelo próprio Maxwell, no mesmo artigo em que a apresentou, por causa das medições dos calores específicos dos gases; ainda assim, seus principais desenvolvimentos vieram depois dessa falsificação.
  \item Um quarto exemplo, tratado à parte por sua importância, é a Revolução Copernicana (ver questão seguinte).
\end{itemize}

\item \textbf{Por que nem falsificacionismo nem indutivismo são compatíveis com a revolução copernicana?}
\begin{itemize}
  \item Em 1543, a teoria de Copérnico enfrentava fortes objeções mecânicas e astronômicas (o argumento da torre, o fato de objetos soltos não caírem da superfície da Terra em rotação, a ausência de paralaxe estelar, o tamanho aparente constante de Vênus e Marte a olho nu), e Copérnico não tinha resposta satisfatória a elas.
  \item Mesmo assim, a teoria não foi abandonada: foi mantida e desenvolvida por cerca de um século e meio, com o trabalho de Galileu, Kepler e Newton, até ser incorporada a uma física abrangente.
  \item Isso contraria o falsificacionismo, que exigiria a rejeição de uma teoria falsificada (ou, na versão sofisticada, sua substituição por uma alternativa já mais falsificável), não sua manutenção obstinada por gerações apesar de falsificações aparentes.
  \item Isso também contraria o indutivismo, pois os novos conceitos de força e inércia não surgiram de observação e experimentação cuidadosas e cumulativas: foram concepções novas e incompletamente formuladas, mantidas por convicção teórica e só gradualmente articuladas.
  \item Somente depois que o novo sistema newtoniano estava plenamente construído, processo que envolveu o trabalho de vários cientistas ao longo de séculos, a nova teoria pôde ser comparada em detalhe com os resultados de observação e experimento.
  \item Conclusão de Chalmers: nenhuma explicação da ciência pode ser aceita como suficiente a menos que acomode episódios históricos como este, o que aponta para explicações mais estruturadas, como os programas de pesquisa de Lakatos e os paradigmas de Kuhn, tratadas nos capítulos seguintes do livro.
\end{itemize}

\end{enumerate}

\bibliographystyle{apalike}
\bibliography{ref}
\end{document}
